% =========================================================================
% ปฏิบัติการที่ 5: การประยุกต์โมเดล Machine Learning บนเบราว์เซอร์
% =========================================================================

\chapter{การประยุกต์โมเดล Machine Learning บนเบราว์เซอร์}
\label{ch:lab05}

\noindent{\large\bfseries\color{cyberblue} การผสานโมเดลจำแนกภาพ Teachable Machine และ TensorFlow.js สู่ฉากเสมือน}

\vspace{0.4cm}

\begin{labobjectivebox}{การประยุกต์โมเดล Machine Learning บนเบราว์เซอร์}
\textbf{วัตถุประสงค์การเรียนรู้}
\begin{enumerate}[leftmargin=1.8cm, itemsep=2pt]
    \item เข้าใจกระบวนการฝึกโมเดลจำแนกภาพถ่ายด้วย Google Teachable Machine
    \item สามารถแปลงและโหลดโมเดล TensorFlow.js เข้าสู่หน่วยความจำเบราว์เซอร์
    \item สามารถเชื่อมโยงผลการพยากรณ์คลาส (Class Prediction) เข้ากับการเปลี่ยนสถานะของวัตถุเสมือน 3 มิติ
\end{enumerate}

\vspace{4pt}
\textbf{เครื่องมือและอุปกรณ์จำเป็น} \enspace Google Teachable Machine (Image Project), TensorFlow.js (@tensorflow/tfjs), ชุดภาพถ่ายตัวอย่าง
\end{labobjectivebox}

\section{หลักการและทฤษฎีพื้นฐาน}

การผสานปัญญาประดิษฐ์เข้ากับเทคโนโลยีเสมือนจริงบนเว็บเบราว์เซอร์ทำได้อย่างมีประสิทธิภาพสูงด้วย TensorFlow.js ซึ่งใช้ประโยชน์จาก WebGL หรือ WebGPU ในการคำนวณแบบขนานบนหน่วยประมวลผลกราฟิก (GPU) ของอุปกรณ์โดยตรง ทำให้ไม่ต้องพึ่งพาเซิร์ฟเวอร์ภายนอก (Zero Server Latency) และรักษาความเป็นส่วนตัวของข้อมูลผู้ใช้งาน

ในปฏิบัติการนี้ ผู้เรียนจะฝึกแบบจำลองโครงข่ายประสาทเทียมแบบคอนโวลูชัน (Convolutional Neural Network: CNN) ที่ผ่านการเรียนรู้ถ่ายโอน (Transfer Learning) จาก MobileNet ผ่านแพลตฟอร์ม Google Teachable Machine เพื่อจำแนกประเภทวัตถุทางวิทยาศาสตร์ เช่น ตัวอย่างแร่ธาตุ ตัวอย่างใบไม้ หรือการแสดงท่าทางพิเศษ ผลลัพธ์จากการพยากรณ์จะอยู่ในรูปของเวกเตอร์ความน่าจะเป็นแบบซอฟต์แมกซ์ (Softmax Probability Vector):
\begin{equation}
P(C_i \mid \mathbf{x}) = \frac{e^{z_i}}{\sum_{j=1}^{K} e^{z_j}}
\end{equation}
เมื่อคลาสใดมีค่าความน่าจะเป็นสูงสุดเกินเกณฑ์ความเชื่อมั่นที่กำหนด ระบบจะเรียกใช้ฟังก์ชันควบคุมโมเดล 3 มิติเพื่อแสดงข้อมูลสารสนเทศ แอนิเมชัน หรือป้ายกำกับโฮโลกราฟิกที่สอดคล้องกัน

\begin{conceptbox}{สาระสำคัญของปฏิบัติการที่ 5}
การเข้าใจโครงสร้างทางคณิตศาสตร์และการประมวลผลเชิงพื้นที่ ช่วยให้นักศึกษาสามารถออกแบบระบบเสมือนจริงที่ตอบสนองต่อผู้ใช้งานได้อย่างเป็นธรรมชาติและแม่นยำสูง ทั้งยังสามารถต่อยอดสู่การพัฒนาแอปพลิเคชันทางวิทยาศาสตร์และอุตสาหกรรมได้อย่างมีประสิทธิภาพ
\end{conceptbox}

\section{ขั้นตอนการปฏิบัติการและการทดลอง}

ให้นักศึกษาดำเนินการสร้างไฟล์โครงการและเขียนคำสั่งตามลำดับขั้นตอนต่อไปนี้ โดยตรวจสอบความถูกต้องของไลบรารีที่เรียกใช้งานและเส้นทางไฟล์

\begin{codebox}{โครงสร้างโค้ดหลักของปฏิบัติการที่ 5}
\begin{lstlisting}[language=HTML]
// การโหลดและรันโมเดล Teachable Machine บนเบราว์เซอร์
const MODEL_URL = './my_model/';
let model, maxPredictions;

async function initAIModel() {
    const modelURL = MODEL_URL + 'model.json';
    const metadataURL = MODEL_URL + 'metadata.json';

    // โหลดโมเดลผ่านไลบรารี tmImage
    model = await tmImage.load(modelURL, metadataURL);
    maxPredictions = model.getTotalClasses();
    console.log(`โมเดลโหลดสำเร็จ: รองรับ ${maxPredictions} คลาส`);
}

async function predictFrame(videoElement) {
    const prediction = await model.predict(videoElement);
    
    for (let i = 0; i < maxPredictions; i++) {
        const className = prediction[i].className;
        const probability = prediction[i].probability.toFixed(2);
        
        // หากตรวจพบวัตถุเป้าหมายด้วยความแม่นยำมากกว่า 85%
        if (probability > 0.85) {
            triggerARResponse(className);
        }
    }
}

function triggerARResponse(detectedClass) {
    // ปรับเปลี่ยนสีและหมุนโมเดลตามคลาสที่ AI ตรวจพบ
    if (detectedClass === "Mineral_Quartz") {
        cube.material.color.setHex(0x38bdf8);
        cube.scale.set(1.5, 1.5, 1.5);
    } else if (detectedClass === "Mineral_Pyrite") {
        cube.material.color.setHex(0xf59e0b);
        cube.scale.set(1.0, 1.0, 1.0);
    }
}
\end{lstlisting}
\end{codebox}

\section{การทดสอบระบบและการบันทึกผล}

เมื่อรันโปรแกรมบนเซิร์ฟเวอร์จำลอง (Local Development Server) ให้ทำการทดสอบการทำงานตามเกณฑ์ประเมินในตารางต่อไปนี้ พร้อมบันทึกผลการสังเกตการณ์ลงในรายงานผลการทดลอง

\begin{table}[htbp]
\centering
\small
\begin{tabularx}{\textwidth}{c X c Y}
\toprule
\textbf{ลำดับ} & \textbf{รายการทดสอบและพฤติกรรมที่คาดหวัง} & \textbf{เกณฑ์ผ่าน} & \textbf{ผลการทดสอบ} \\
\midrule
1 & การเริ่มต้นระบบและการเข้าถึงสิทธิ์ฮาร์ดแวร์ & ภายใน 3 วินาที & ผ่าน / ไม่ผ่าน \\
2 & ความเสถียรของการตรวจจับและการคำนวณพิกัด & อัตราความแม่นยำ $> 90\%$ & ผ่าน / ไม่ผ่าน \\
3 & อัตราการแสดงผลกราฟิกต่อเนื่อง (Framerate) & ไม่ต่ำกว่า 55 FPS & ผ่าน / ไม่ผ่าน \\
4 & การตอบสนองต่อปฏิสัมพันธ์เชิงพื้นที่ & ความหน่วง $< 40$ ms & ผ่าน / ไม่ผ่าน \\
\bottomrule
\end{tabularx}
\caption{ตารางประเมินผลการทำงานของระบบในปฏิบัติการที่ 5}
\label{tab:eval_lab05}
\end{table}

\section{ภารกิจท้าทายและการต่อยอด}

\begin{activitybox}{กิจกรรมส่งเสริมทักษะขั้นสูง}
ให้นักศึกษาฝึกสอนโมเดล Teachable Machine ด้วยภาพถ่ายการแสดงสัญลักษณ์มือ 3 แบบ (OK, Peace, Thumbs Up) และเขียนเงื่อนไขเพื่อเปลี่ยนรูปแบบแสงไฟและระบบเสียงสังเคราะห์ในฉาก WebAR ตามสัญลักษณ์ที่ตรวจจับได้
\end{activitybox}

\section{คำถามท้ายการทดลอง}

ให้นักศึกษาตอบคำถามเชิงวิเคราะห์ต่อไปนี้ลงในสมุดบันทึกปฏิบัติการ

\begin{enumerate}[leftmargin=1.8cm, itemsep=6pt]
    \item จงอธิบายข้อดีของการใช้ Transfer Learning บน MobileNet สำหรับการประมวลผลบนเบราว์เซอร์
    \item เหตุใดการรันโมเดล AI ผ่าน WebGL Shader จึงมีความเร็วสูงกว่าการคำนวณผ่าน JavaScript เอนจินแบบเดิม
    \item หากภาพวิดีโอมีพื้นหลังซับซ้อน จะส่งผลต่อค่าความน่าจะเป็นในสมการ Softmax อย่างไร และมีวิธีแก้ปัญหาในขั้นตอนการเตรียมข้อมูลอย่างไร
\end{enumerate}

